\documentclass{article}
\usepackage{../../Self_Style}

\title{Latex Template}
\author{Zih-Yu Hsieh}
\date{\today}

\begin{document}
\maketitle

\section*{1}

\begin{ques}\label{q1}
\textbf{You just know the answer}

A common problem in QM is ``find the energies of the following Hamiltonian.'' The idea
is to do as little math as possible but to build intuition. Describe the lowest several
eigenenergies of the following Hamiltonians. It is up to you to decide what value of ``several''
best illustrates what is going on. Answers are short.

\textbf{Example:} What are the energies of the two separate simple harmonic oscillators, each
with frequency $\omega$?

\textbf{Answer:} The oscillators have energies
\[
E_1 = \hbar\omega(1/2 + n_1), \quad n_1 = [0,1,2,3,\ldots]
\]
\[
E_2 = \hbar\omega(1/2 + n_2), \quad n_2 = [0,1,2,3,\ldots]
\]
so the combined oscillators have energy
\[
E = E_1 + E_2 = \hbar\omega(1 + n_1 + n_2)
= \hbar\omega[1,2,2,3,3,3,4,4,4,4,\ldots]
\]

\begin{enumerate}
\item Particle of mass $m$ in one dimensional box of width $a$.
\item Particle of mass $m$ in two dimensional box of width $a$ and length $b$.
\item Two disconnected particles, each of mass $m$, one in a 1D box of width $a$ and the
other in a 1D box of width $b$.
\item Particle of mass $m$ confined to a circular ring of circumference $a$.
\item Particle of mass $m$ in a one dimensional potential $V = \frac{1}{2}m\omega^2 x^2$.
\item Paricle of mass $m$ in a two dimensional potential
$V = \frac{1}{2}m\omega^2(x^2 + y^2)$.
\end{enumerate}
\end{ques}

\begin{proof}
\end{proof}

\newpage

\section*{2}

\begin{ques}\label{q2}
\textbf{Two-body to one-body} (Griffiths 5.1 in both editions)

Given my history, I will not transcribe these anymore unless I add something to it, but feel
free to email if you need me to.
\end{ques}

\begin{proof}
\end{proof}

\newpage

\section*{3}

\begin{ques}\label{q3}
\textbf{SHO in 3D} (Modified Griffiths 3rd ed.\ 4.46, or 2nd ed.\ 4.39)

Consider a particle of mass $m$ for which the potential is
\[
V(r) = \frac{1}{2} m\omega^2 r^2
\]

\begin{enumerate}
\item[(a)] Using separation of variables in Cartesian coordinates, show that this factorizes into
a sum of three 1D harmonic oscillators, and use your knowledge of 1D SHO to
determine the allowed energies.
\item[(b)] Determine the degeneracy $d(n)$ (i.e.\ the number of states with the same energy as
the $n$th energy eigenvale $E_n$.
\item[(c)] As you hopefully found in part (b), the ground state $\psi_0(x,y,z)$ is non-degenerate.
However the first excited energy eigenvalue $E_1$ is threefold degenerate, with eigen-
states:
\[
\psi_{1,0,0}(x,y,z) = \tilde{\psi}_1(x)\tilde{\psi}_0(y)\tilde{\psi}_0(z)
\]
\[
\psi_{0,1,0}(x,y,z) = \tilde{\psi}_0(x)\tilde{\psi}_1(y)\tilde{\psi}_0(z)
\]
\[
\psi_{0,0,1}(x,y,z) = \tilde{\psi}_0(x)\tilde{\psi}_0(y)\tilde{\psi}_1(z)
\]
where $\tilde{\psi}_i$ denotes an eigenstate of the 1D harmonic oscillator. Are any of the linear
combinations
\[
\frac{1}{\sqrt{3}}(\psi_{1,0,0} + \psi_{0,1,0} + \psi_{0,0,1}),
\quad
\frac{1}{\sqrt{2}}(\psi_{1,0,0} + \psi_{0,1,0}),
\quad
\frac{1}{\sqrt{2}}(\psi_{1,0,0} + i\psi_{0,0,1})
\]
also energy eigenstates? In each case, why or why not?
\end{enumerate}
\end{ques}

\begin{proof}
\end{proof}

\newpage

\section*{4}

\begin{ques}\label{q4}
\textbf{1D to 3D} (Griffiths 4.1, same in both editions)
\end{ques}

\begin{proof}
\end{proof}

\newpage

\section*{5}

\begin{ques}\label{q5}
\textbf{Cubical well} (Griffiths 4.2, same in both editions)
\end{ques}

\begin{proof}
\end{proof}

\newpage

\section*{6}

\begin{ques}\label{q6}
\textbf{Practice with Spherical Harmonics} (Griffiths 4.4 in 3rd edition, 4.3 in 2nd edition)
\end{ques}

\begin{proof}
\end{proof}

\end{document}
